\documentclass[10pt,a4paper]{amsart}
\usepackage[margin=19mm]{geometry}
\usepackage{amsmath,amssymb,mathtools}
\usepackage[scr=rsfs,cal=euler]{mathalfa}
\usepackage{xfrac}
\usepackage[unicode,hidelinks,bookmarks=false]{hyperref}
\newcommand{\ie}{i.e.\ }
\newcommand{\cf}{cf.\ }
\newcommand{\resp}{respectively\ }
\newcommand{\Subsection}{\subsection}
\providecommand{\nobreakdash}{\nobreak\hbox{-}}
\setlength{\emergencystretch}{2em}
\multlinegap0pt
\usepackage{amssymb}
\usepackage{xcolor}
\usepackage{enumerate}
\newtheorem{defn}{Definition}
\newtheorem{thm}{Theorem}
\title{On the subdirect product of graph bundles}
\author{Yanga Bavuma, Francesco G. Russo, Stefano Spessato}
\date{Source: arXiv:2507.14530v1 (CC BY 4.0)}
\begin{document}
\maketitle

\noindent\emph{Source and licence.} On the subdirect product of graph bundles. Version arXiv:2507.14530v1 (CC BY 4.0), distributed by arXiv under the Creative Commons Attribution 4.0 International licence, http://creativecommons.org/licenses/by/4.0/, as recorded in the arXiv licence metadata for this exact version and shown on the abstract page https://arxiv.org/abs/2507.14530v1. Full licence notice: \texttt{licenses/arxiv-2507.14530-CC-BY-4.0.txt}.\par\smallskip

\noindent\textit{Editorial scope.} Counted unit: the article's thm sumstab (source0.tex lines 1050--1055). The article's complete original proof of it is delivered with it (source0.tex lines 1056--1080, one \texttt{proof} environment of 25 lines). No mathematics is written, repaired or re-derived: the statement and the proof are the article's own lines, in the article's own notation, and no retained line is substituted or re-flowed.  Retained uncounted as the source-local context the proof needs: lines 585--597.  Nothing was omitted from any retained span, so the package carries no omission record.  No retained line contains a citation, so the wrapper carries no bibliography and no bibliography entry.  No retained line contains a figure environment, an image-inclusion command or drawing code, so no image is delivered and the package declares no asset dependency.  Editorial formatting only, in addition: an AMS \texttt{amsart} preamble that restates the article's own declarations and macros used by the retained lines, namely the environments \texttt{defn}, \texttt{thm} on separate counters, together with the standard packages those lines need.  The retained environments print in this document as Definition 1, Theorem 1 in order of appearance, whereas the article prints the same items with the numbering of its own full document.  The complete native source of the article as distributed in the arXiv e-print for this exact version is bundled unchanged as \texttt{sources/181857/source0.tex}.
\begin{defn}\label{pullback}
The \textbf{pullback along } $f$ \textbf{of } $(X, p, \Gamma')$ is the $F$-graph bundle $(f^*X, f^*p, \Gamma)$ where $f^*X$ is the graph whose vertex set is 
\begin{equation*}
    V_{f^* X} = \{(v, x) \in V_{\Gamma} \times V_{X} \mid f(v) = p(x)\}
\end{equation*}
and $(v,x) \sim (w, y)$ if and only if one of the following three conditions hold
\begin{itemize}
    \item[\rm (1).] $v = w$, and $x\sim y$ (type I edges),
    \item[\rm (2).] $v\sim w$, $f(v) = f(w)$ and $x = y$ (type II edges),
    \item[\rm (3).] $v\sim w$, $f(v) \sim f(w)$ and $x \sim y$ (type III edges).
\end{itemize}
Finally $f^*p: f^*X \longrightarrow \Gamma$ is the morphism such that for each vertex $(v,x)$ in $f^*X$ is defined as $f^*p(v,x) = v$. We denote this graph bundle by $f^*(X, p, \Gamma')$.
\end{defn}

\begin{thm}\label{sumstab}
Let $(X_i, p_i, \Gamma)$ be two $F_i$-graph bundles over the same graph $\Gamma$ and  $f: \Gamma' \longrightarrow \Gamma$  a morphism of graphs. Then
\begin{equation}
f^*(X_1 \boxplus X_2, (p_1,p_2), \Gamma) \cong f^*(X_1, p_1, \Gamma) \boxplus f^*(X_2, p_2, \Gamma).
\end{equation}
\end{thm}

\begin{proof}
The vertex set of $f^*(X_1 \boxplus X_2)$ is
\begin{equation}
\{(v,x, y) \in V_{\Gamma'}\times V_{X_1} \times V_{X_2} \mid f(v)= p_1(x)= p_2(y)\}.    
\end{equation}
On the other hand, the vertex set of $f^*X_1 \boxplus f^*X_2$ is
\begin{equation}
\{(v,x, v, y) \in V_{\Gamma'} \times V_{X_1} \times V_{\Gamma'} \times V_{X_2} \mid f(v)= p_1(x)= p_2(y)\}.    
\end{equation}
Let $\Phi$ be the map defined for each vertex in $f^*(X_1 \boxplus X_2)$ by \begin{equation}\Phi(v,x, y) = (v,x, v, y).\end{equation} Clearly $\Phi$ commutes with the projections: so, if it is an isomorphism of graphs, then it is also an isomorphism of graph bundles. Since $f^*(X_1  \ \boxplus \  X_2, (p_1,p_2), \Gamma)$ and $f^*(X_1, p_1, \Gamma) \boxplus f^*(X_2, p_2, \Gamma)$ are both $F_1  \ \square  \ F_2$-graph bundle on the same base $\Gamma$, it is sufficient to show that $\Phi$ is a morphism of graph which preserves the edges, indeed two graph bundles with the same fiber and the same base also have the same number of edges. This is a consequence of the fact that a graph bundle is locally a Cartesian product. Let's argue according to the types of the edge as per Definition \ref{pullback}.

\begin{enumerate}
    \item[(1).] 
Let $(v,x, y)\sim(v,x',y')$. These vertices form an edge of type $I$ for $f^*(X_1 \boxplus X_2)$. In particular either $x \sim x'$ and $y=y'$, or $x = x'$ and $y \sim y'$. In the first case $\Phi(v,x, y) = (v,x, v, y)\sim(v,x',v, y')= \Phi(v,x', y')$ is an edge of type II for $f^*(X_1, p_1, \Gamma)  \ \boxplus \ f^*(X_2, p_2, \Gamma)$. In the second case, it is an edge of type I.


    \item[(2).] 
Let $(v,x, y)\sim(w,x,y)$ be an edge of type II. Then $\Phi(v,x, y) = (v,x, v, y)\sim(w,x,w, y)= \Phi(w,x, y)$ is also an edge of type II for $f^*(X_1, p_1, \Gamma)  \ \boxplus  \ f^*(X_2, p_2, \Gamma)$.


    \item[(3).] 
Let $(v,x, y)\sim(w,x',y')$ be an edge of type III. Then we have $f(v) \sim f(w)$ and $(x,y) \sim (x',y')$.  Since $(p_1,p_2)(x,y)= f(v)$ and $(p_1,p_2)(x',y')= f(w)$, we find that $x \sim x'$ and $y\sim y'$. So the vertices $(v,x) \sim (w,x')$ in $f^*X_1$, $(v,y) \sim (w,y')$ in $f^*X_2$ and $\Phi(v,x, y) = (v,x, v, y)\sim(w,x',w, y')= \Phi(w,x', y')$ is an edge of type III in $f^*(X_1, p_1, \Gamma)  \ \boxplus  \ f^*(X_2, p_2, \Gamma)$. 
\end{enumerate}
The result follows.
\end{proof}

\end{document}
